% =========================================================
% Stage 0 S0-C
% =========================================================
\chapter{S0-C --- Complexity \& Numeric Safety}

% =========================================================
% Part A
% =========================================================
\section{Part A --- 위치와 선수지식}

\subsection{현재 위치}
S0-A에서 C++의 기본 실행 구조와 scalar 상태 전이를, S0-B에서 \code{string}, \code{vector}, \code{pair}, iterator, \code{sort}, \code{reverse}, reference와 \code{const}를 다뤘다면, S0-C에서는 작성한 코드가 \textbf{주어진 입력 크기에서 실제로 실행 가능한지}, 그리고 \textbf{정수 연산이 실제 C++ 평가 순서에서 안전한지}를 체계적으로 판단한다.

이번 Learning Unit의 핵심 범위는 다음과 같다.
\begin{itemize}
  \item Big-O와 대표 성장률
  \item 입력 크기에서 허용 가능한 복잡도 추정
  \item 순차 반복, 중첩 반복, 정렬의 비용
  \item 원소 개수 $N$과 총 payload 크기 $L$, 그룹 전체 원소 수 $T$의 구분
  \item total space와 auxiliary space의 구분
  \item \code{int}, \code{long long}, container size type의 범위와 역할
  \item 정수 overflow와 intermediate expression
  \item \code{1LL}, 명시적 cast, usual arithmetic conversion
  \item 전역 upper bound와 실제 C++ expression type을 함께 검증하는 numeric-safety proof
\end{itemize}

\subsection{왜 지금 배우는가}
Competitive Programming에서 정답 알고리즘은 단순히 ``논리적으로 맞는 방법''이 아니다. 같은 논리를 구현해도 $O(N^2)$ 알고리즘이 제한시간을 넘거나, 최종 결과는 64-bit에 들어가지만 중간의 \code{int * int}에서 먼저 overflow가 발생하면 오답이 된다.

또한 문자열이나 여러 그룹이 주어지는 문제에서는 입력 크기를 단순히 ``원소 수 $N$'' 하나로만 표현할 수 없다. 예를 들어 $N$개의 문자열 전체 길이가 $L$이라면, 모든 문자열을 출력하거나 뒤집는 비용은 $O(N)$이 아니라 실제 처리 문자 수에 비례하는 $O(L)$이 된다.

\begin{keybox}[S0-C의 핵심 질문]
코드를 작성한 뒤 다음 네 질문에 답할 수 있어야 한다.
\begin{enumerate}
  \item 실제 input size는 무엇인가?
  \item 그 input size에 대해 몇 번의 연산을 수행하는가?
  \item 코드의 각 중요한 중간값은 최대 얼마나 커질 수 있는가?
  \item 그 중간값은 실제 C++에서 어떤 타입으로 계산되는가?
\end{enumerate}
\end{keybox}

\subsection{Core Decision Boundary}
S0-C에서는 다음 판단 경계를 핵심으로 본다.
\begin{center}
\begin{tabularx}{\textwidth}{p{4.3cm}X}
\toprule
판단 상황 & 핵심 질문 \\
\midrule
$O(N)$ vs $O(N\log N)$ vs $O(N^2)$ & 입력 크기에서 어느 성장률까지 현실적으로 허용되는가? \\
$N$ vs $L$ vs $T$ & 실제로 반복하거나 복사하거나 출력하는 데이터의 총량은 무엇인가? \\
\code{int} vs \code{long long} & 최종값뿐 아니라 각 intermediate가 어느 범위까지 커지는가? \\
수학적 식 vs C++ 식 & 같은 식처럼 보여도 C++에서 어떤 subexpression이 먼저 어떤 타입으로 평가되는가? \\
정확한 최대값 vs 보수적 upper bound & 정확한 최대값이 필요한가, 아니면 모든 입력을 덮는 안전한 상계면 충분한가? \\
\bottomrule
\end{tabularx}
\end{center}

\subsection{Immediate Coverage Floor}
이 Unit을 즉시 통과하기 위한 최소 증거는 다음과 같이 본다.
\begin{itemize}
  \item 입력 크기와 반복 구조로부터 시간복잡도를 독립적으로 설명한다.
  \item total space와 auxiliary space를 구분한다.
  \item 정수 범위에서 최종값과 주요 intermediate를 구분해 upper bound를 제시한다.
  \item 실제 C++ operand type과 evaluation order를 확인하여 overflow 안전성을 설명한다.
\end{itemize}

% =========================================================
% Part B
% =========================================================
\section{Part B --- 학습 목표}
이 Unit을 마치면 다음 행동을 독립적으로 수행할 수 있어야 한다.

\begin{enumerate}
  \item $O(1)$, $O(\log N)$, $O(N)$, $O(N\log N)$, $O(N^2)$의 성장 차이를 설명한다.
  \item 입력 제한을 보고 명백히 불가능한 알고리즘을 배제한다.
  \item 순차 반복은 합하고, 중첩 반복은 반복 횟수를 곱하거나 총 iteration 수를 직접 세어 복잡도를 계산한다.
  \item 문자열 문제에서 원소 개수 $N$과 총 문자 수 $L$을 구분한다.
  \item 그룹 입력에서 그룹 수 $N$과 총 원소 수 $T$를 구분한다.
  \item \code{sort}, \code{reverse}, 문자열 복사, 문자열 출력의 비용을 실제 처리 범위 크기로 설명한다.
  \item total space와 auxiliary space를 별도로 설명한다.
  \item \code{int}와 \code{long long}의 전형적인 범위를 기억하고 입력 제한에서 필요한 타입을 선택한다.
  \item \code{long long x = a * b;}가 항상 안전하지 않은 이유를 설명한다.
  \item \code{1LL} 또는 명시적 cast를 이용해 multiplication chain을 원하는 타입으로 승격시킨다.
  \item 최종 답뿐 아니라 raw intermediate product의 upper bound를 계산한다.
  \item 한 개의 가능한 입력 예시와 모든 입력을 덮는 전역 upper bound를 구분한다.
  \item \code{string::size()}와 \code{vector::size()}의 반환형이 signed \code{int}가 아님을 인식한다.
  \item numeric-safety proof를 ``constraint $\to$ bound $\to$ actual type $\to$ capacity'' 순서로 작성한다.
\end{enumerate}

% =========================================================
% Part C
% =========================================================
\section{Part C --- 개념과 원리}

\subsection{C1. Big-O는 입력 크기가 커질 때의 성장률이다}
시간복잡도는 특정 컴퓨터에서 몇 초가 걸리는지를 직접 표현하지 않는다. 입력 크기 $N$이 커질 때 필요한 연산 수가 어떤 속도로 증가하는지를 표현한다.

대표적인 성장률은 다음과 같다.
\[
O(1),\quad O(\log N),\quad O(N),\quad O(N\log N),\quad O(N^2).
\]

예를 들어 한 번의 선형 순회는 보통 $O(N)$이고, 정렬은 일반적으로 $O(N\log N)$이다. 모든 ordered pair를 직접 확인하는 이중 반복문은 흔히 $O(N^2)$이 된다.

\subsection{C2. $O$, $\Omega$, $\Theta$의 역할}
Big-O는 성장률의 상계를 표현한다. 반대로 $\Omega$는 하계를, $\Theta$는 상계와 하계가 같은 차수임을 표현할 때 사용한다.

CP에서 ``이 알고리즘은 $O(N)$이다''라고 말할 때는 보통 제한시간을 판단하기 위한 upper-bound 관점이 가장 중요하다. 다만 실제 성장률이 선형임을 더 정확하게 표현하고 싶다면 $\Theta(N)$이라고 쓸 수 있다.

\subsection{C3. 순차 반복과 중첩 반복}
두 개의 선형 반복문이 순서대로 실행되면
\[
O(N)+O(N)=O(N)
\]
이다.

반면 각 $i$마다 다시 $N$개를 순회하면
\[
O(N)\times O(N)=O(N^2)
\]
가 된다.

다만 ``for문이 두 개''라는 외형만 보고 무조건 $O(N^2)$라고 판단하면 안 된다. 그룹마다 $M_g$번 반복하고
\[
T=\sum_g M_g
\]
라면 내부 반복 전체의 총 iteration 수는 $T$이므로 $O(T)$일 수 있다.

\subsection{C4. 입력 제한에서 실행 가능성을 역산하기}
문제의 제한은 알고리즘 후보를 걸러내는 정보다. $N$이 매우 큰데 모든 pair를 확인하는 $O(N^2)$ 접근을 사용하면 연산 수가 급격히 증가한다. 반대로 작은 입력에서는 더 단순한 brute force가 충분할 수 있다.

핵심 절차는 다음과 같다.
\begin{enumerate}
  \item 실제 입력 크기 변수를 찾는다.
  \item 후보 알고리즘의 iteration 수를 그 변수로 표현한다.
  \item 더 빠른 접근이 필요한지 판단한다.
\end{enumerate}

\subsection{C5. 원소 개수 $N$과 총 payload $L$을 구분하기}
$N$개의 문자열이 있고 각 문자열 길이를 $|s_i|$라 하자. 전체 문자 수를
\[
L=\sum_{i=0}^{N-1}|s_i|
\]
라고 정의하면, 모든 문자열을 한 번씩 뒤집는 총 비용은
\[
O(|s_0|)+\cdots+O(|s_{N-1}|)=O(L)
\]
이다.

모든 문자열을 출력하는 비용도 총 $L$개의 문자를 써야 하므로 $O(L)$이다. ``출력문이 $N$번 실행된다''는 사실만으로 $O(N)$이라고 결론내리면 실제 payload 비용을 놓칠 수 있다.

\subsection{C6. 그룹 수 $N$과 총 원소 수 $T$}
각 그룹 $g$의 크기가 $M_g$이고
\[
T=\sum_{g=0}^{N-1}M_g
\]
라면,
\begin{lstlisting}
for (int g = 0; g < N; ++g) {
    for (int j = 0; j < M[g]; ++j) {
        // O(1)
    }
}
\end{lstlisting}
의 내부 body가 실행되는 횟수는 $N^2$이 아니라 정확히 $T$다. 따라서 전체가 $O(N+T)$ 또는 조건에 따라 $O(T)$가 될 수 있다.

\subsection{C7. 자주 쓰는 STL 연산의 비용}
\begin{center}
\begin{tabularx}{\linewidth}{p{5.2cm}X}
\toprule
연산 & 대표적인 비용 \\
\midrule
\code{v[i]}, \code{v.size()} & $O(1)$ \\
\code{push\_back} & amortized $O(1)$ \\
전체 순회 & $O(N)$ \\
\code{reverse(first,last)} & 뒤집는 범위 길이에 비례 \\
\code{sort(first,last)} & 일반적으로 $O(N\log N)$ \\
\code{string} 복사 & 문자열 길이에 비례 \\
문자열 출력 & 출력하는 문자 수에 비례 \\
\bottomrule
\end{tabularx}
\end{center}

\subsection{C8. total space와 auxiliary space}
입력을 저장하는 공간과 알고리즘이 추가로 사용하는 공간은 구분해서 설명할 수 있다.

예를 들어 길이 $N$의 vector를 저장하고 scalar 변수 몇 개만 추가로 사용한다면
\begin{itemize}
  \item total space: $O(N)$
  \item auxiliary space: $O(1)$
\end{itemize}
로 설명할 수 있다.

반대로 입력을 streaming으로 읽어 전부 저장하지 않는다면 total space 자체가 $O(1)$일 수도 있다.

\subsection{C9. \code{int}와 \code{long long}의 범위}
CP 환경에서 보통 다음 정도의 크기를 기준으로 생각한다.
\begin{itemize}
  \item signed \code{int}: 약 $-2.1\times10^9$에서 $2.1\times10^9$
  \item signed \code{long long}: 약 $-9.22\times10^{18}$에서 $9.22\times10^{18}$
\end{itemize}

중요한 것은 변수 하나의 입력값이 아니라 \textbf{계산 과정에서 만들어질 최대값}이다. 입력 하나가 $10^9$ 이하라도 $10^9\times10^9$은 \code{int} 범위를 훨씬 넘는다.

\subsection{C10. destination type과 expression type은 다르다}
다음 코드는 안전하지 않다.
\begin{lstlisting}
int a = 100000;
int b = 100000;
long long x = a * b;
\end{lstlisting}

\code{x}가 \code{long long}인 것은 \textbf{대입받는 변수의 타입}이다. 오른쪽의 \code{a * b}는 두 operand가 \code{int}이므로 먼저 \code{int} 연산으로 평가된다. 결과가 \code{long long}에 대입되기 전에 overflow할 수 있다.

\begin{keybox}[기억할 문장]
\code{long long} 변수에 저장한다고 해서 그 변수에 도달하기 전의 모든 연산이 자동으로 \code{long long}이 되는 것은 아니다.
\end{keybox}

\subsection{C11. early promotion: \code{1LL}과 cast}
곱셈 체인을 처음부터 \code{long long}으로 만들려면 다음처럼 쓸 수 있다.
\begin{lstlisting}
long long x = 1LL * a * b;
\end{lstlisting}

또는
\begin{lstlisting}
long long x = static_cast<long long>(a) * b;
\end{lstlisting}
처럼 명시적으로 승격할 수 있다.

cast의 위치도 중요하다.
\begin{lstlisting}
long long x = static_cast<long long>(a + b) * c;
\end{lstlisting}
에서는 \code{a + b}가 먼저 \code{int}로 계산된다. 그 합 자체가 overflow할 수 있다면 다음처럼 더 일찍 승격해야 한다.
\begin{lstlisting}
long long x = (static_cast<long long>(a) + b) * c;
\end{lstlisting}

\subsection{C12. 최종값뿐 아니라 cumulative intermediate를 확인한다}
다음 식을 생각하자.
\begin{lstlisting}
long long x = 1LL * n * (n + 1) * value / 2;
\end{lstlisting}

최종적으로 2로 나눈 값이 \code{long long} 범위 안에 들어가더라도, 나누기 전의 raw product가 더 클 수 있다. C++은 왼쪽부터 multiplication chain을 평가하므로 각 cumulative subexpression의 최대 크기를 확인해야 한다.

필요하면 먼저 2로 나누어 intermediate를 줄일 수 있다.
\begin{lstlisting}
long long a = n;
long long b = n + 1;
if (a % 2 == 0) a /= 2;
else b /= 2;
long long x = a * b * value;
\end{lstlisting}

\subsection{C13. 가능한 한 입력 예시와 upper bound는 다르다}
``이 입력에서는 값이 작다''는 것은 전역 안전성 증명이 아니다. 모든 valid input을 덮는 상계를 보여야 한다.

예를 들어 실제 가능한 최대값이 $2.5\times10^{13}$인데 $5\times10^{12}$라고 적으면, 이것은 단순한 근사 오차가 아니라 upper bound가 아니다. 반대로 $3\times10^{13}$이라면 실제 최대보다 크므로 느슨하지만 유효한 upper bound가 될 수 있다.

\subsection{C14. upper bound 산술을 안전하게 계산하는 방법}
다음 표현식의 bound를 계산한다고 하자.
\[
j^2A,\qquad j<5\times10^4,\quad A\le10^4.
\]

계수와 $10$의 지수를 분리하면
\[
(5\times10^4)^2(10^4)
=25\times10^{12}
=2.5\times10^{13}.
\]

큰 수를 암산할 때는 마지막 정규화까지 적으면 지수 실수를 줄일 수 있다.

\subsection{C15. \code{size()}의 반환형과 signed/unsigned}
\code{string::size()}와 \code{vector::size()}는 일반적인 \code{int}가 아니라 각 container의 \code{size\_type}을 반환한다. 이 타입은 보통 unsigned 계열이다.

따라서 다음과 같은 식은 operand type을 실제로 확인해야 한다.
\begin{lstlisting}
1LL * (i + 1) * s.size()
\end{lstlisting}

선두의 \code{1LL} 때문에 앞부분은 \code{long long}으로 시작하지만, 뒤에서 unsigned \code{size\_type}과 결합될 때 usual arithmetic conversion이 어떻게 적용되는지 생각해야 한다. 이 문제처럼 모든 값이 음이 아니고 매우 작은 범위라면 결과가 안전할 수 있지만, ``\code{1LL}이 있으므로 뒤의 모든 식은 무조건 signed \code{long long}''이라고 기계적으로 단정해서는 안 된다.

\subsection{C16. Numeric-safety proof의 표준 순서}
S0-C에서는 다음 순서를 기본 proof template으로 사용한다.
\begin{enumerate}
  \item \textbf{Constraint}: 각 변수의 범위를 적는다.
  \item \textbf{Global bound}: 최종 answer의 전역 upper bound를 적는다.
  \item \textbf{Intermediate bounds}: 실제 코드에서 overflow 위험이 있는 cumulative subexpression의 상계를 적는다.
  \item \textbf{Actual C++ type}: 그 subexpression이 실제로 어떤 타입으로 평가되는지 적는다.
  \item \textbf{Capacity comparison}: bound와 타입의 최대 범위를 비교해 Safe/Unsafe를 결론낸다.
\end{enumerate}

\begin{checklistbox}[짧은 numeric-safety 체크]
\begin{itemize}
  \item[$\square$] 실제 최대 입력을 모두 덮는 upper bound인가?
  \item[$\square$] 나누기 전 raw product를 확인했는가?
  \item[$\square$] destination type만 보고 있지 않은가?
  \item[$\square$] cast가 너무 늦게 적용되지 않는가?
  \item[$\square$] \code{size\_type} 같은 unsigned operand가 섞이지 않는가?
\end{itemize}
\end{checklistbox}

% =========================================================
% Part D
% =========================================================
\section{Part D --- Worked Example}

\subsection{D1. 모든 pair product의 합}
길이 $N$의 음이 아닌 정수 수열 $A$가 있을 때
\[
S=\sum_{0\le i<j<N}A_iA_j
\]
를 계산한다고 하자.

조건은 다음과 같다.
\begin{TextBlock}
1 <= N <= 200000
0 <= A_i <= 10000
\end{TextBlock}

모든 pair를 직접 확인하면 $O(N^2)$이다. 대신 $j$번째 값을 읽을 때 이전 값의 합을 \code{prefix}라고 두면, $A_j$가 만드는 새 pair의 합은
\[
A_j(A_0+\cdots+A_{j-1})
\]
이다.

\begin{lstlisting}
#include <iostream>
using namespace std;

int main() {
    int N;
    cin >> N;

    long long prefix = 0;
    long long answer = 0;

    for (int i = 0; i < N; ++i) {
        long long x;
        cin >> x;
        answer += prefix * x;
        prefix += x;
    }

    cout << answer << '\n';
}
\end{lstlisting}

\subsection{D2. 시간복잡도와 공간복잡도}
각 원소를 한 번만 읽으므로 시간복잡도는 $O(N)$이다. 전체 입력을 저장하지 않고 \code{prefix}, \code{answer}, \code{x}만 유지하므로 total space와 auxiliary space 모두 $O(1)$이다.

\subsection{D3. Numeric safety}
최종값은 모든 원소가 $10000$일 때 가장 크게 잡을 수 있다.
\[
S\le \binom{200000}{2}\cdot10^8
=1.99999\times10^{18},
\]
따라서 signed \code{long long} 범위 안이다.

또한 \code{prefix * x}에서
\[
\code{prefix}\le 199999\cdot10000<2\times10^9,
\]
\[
\code{prefix}\cdot x<2\times10^{13},
\]
이고 두 operand가 모두 \code{long long}이므로 곱셈 자체도 signed 64-bit로 안전하게 평가된다.

\begin{keybox}[Worked Example의 핵심]
알고리즘을 $O(N^2)$에서 $O(N)$으로 바꾸는 것과, 그 $O(N)$ 코드가 정수 범위에서도 안전한지 증명하는 것은 별개의 단계다. S0-C에서는 두 단계를 모두 완료해야 한다.
\end{keybox}

% =========================================================
% Part E
% =========================================================
\section{Part E --- 중간평가 문제와 보강 학습}

\subsection{E-1. Total Pair Gap}
\begin{problembox}{중간평가 문제 E-1 --- Total Pair Gap}
비내림차순으로 정렬된 정수 수열 $A_0,\dots,A_{N-1}$이 주어진다. 다음 값을 계산하라.
\[
S=\sum_{0\le i<j<N}(A_j-A_i).
\]

\textbf{구현 및 제출 조건}
\begin{itemize}
  \item C++17 또는 C++20을 사용한다.
  \item 목표 시간복잡도는 $O(N)$이다.
  \item 목표 auxiliary space는 $O(1)$이다.
  \item numeric-safety 설명에서 최종값이 signed 64-bit에 들어감을 전역적으로 증명한다.
  \item 완전한 코드, 시간복잡도, 공간복잡도, numeric-safety proof, edge case를 제출한다.
\end{itemize}

\textbf{제약 조건}
\begin{TextBlock}
1 <= N <= 100000
0 <= A_i <= 1000000000
A is sorted in nondecreasing order.
\end{TextBlock}
\end{problembox}

\begin{answerbox}[학습자 제출 답안]
\begin{lstlisting}
#include <iostream>
using namespace std;

int main() {
    int N = 0;
    cin >> N;
    long long s = 0;
    for (int i = 0; i <= N - 1; i += 1) {
        long long A = 0;
        cin >> A;
        s += (2 * i - N + 1) * A;
    }
    cout << s;
}
\end{lstlisting}

제출한 식은 각 $A_i$가 오른쪽 pair에서는 $i$번 더해지고 왼쪽 pair에서는 $N-1-i$번 빠지는 것을 이용한 contribution formula다.
\end{answerbox}

\subsection{보강: 한 feasible input은 global bound가 아니다}
한 입력 예시에서 값이 작다는 사실만으로 전체 입력 범위의 안전성을 증명할 수 없다. 이 문제에서 값 $0$과 $10^9$를 적절히 나누면 cross pair가 가장 많이 만들어질 수 있고, 보수적으로
\[
S\le \left\lfloor\frac{N^2}{4}\right\rfloor10^9
\le2.5\times10^{18}
\]
로 잡을 수 있다. 중요한 것은 정확한 최대 구조를 외우는 것이 아니라, ``모든 valid input을 덮는 상계인가?''를 확인하는 것이다.

\subsection{E-2. Sum of All Subarray Sums}
\begin{problembox}{중간평가 재평가 문제 E-2 --- Sum of All Subarray Sums}
길이 $N$의 음이 아닌 정수 수열 $A$가 주어진다. 모든 연속 부분배열의 원소 합을 다시 모두 더한 값을 출력하라.

\textbf{구현 및 제출 조건}
\begin{itemize}
  \item C++17 또는 C++20을 사용한다.
  \item 목표 시간복잡도는 $O(N)$이다.
  \item 목표 auxiliary space는 $O(1)$이다.
  \item 최종 answer와 실제 코드의 intermediate expression이 signed 64-bit에서 안전한지 설명한다.
\end{itemize}

\textbf{제약 조건}
\begin{TextBlock}
1 <= N <= 100000
0 <= A_i <= 1000
\end{TextBlock}
\end{problembox}

\begin{answerbox}[학습자 제출 코드의 핵심]
\begin{lstlisting}
for (int i = 0; i <= N - 1; i += 1) {
    long long A = 0;
    cin >> A;
    s += (i + 1) * (N - i) * A;
}
\end{lstlisting}

각 $A_i$는 시작점 선택이 $i+1$개, 끝점 선택이 $N-i$개이므로 $(i+1)(N-i)$개의 부분배열에 포함된다.
\end{answerbox}

\subsection{보강: destination type보다 먼저 발생하는 overflow}
위 코드에서 \code{i}와 \code{N}이 \code{int}이면
\begin{lstlisting}
(i + 1) * (N - i)
\end{lstlisting}
가 먼저 \code{int * int}로 계산된다. $N=100000$ 부근에서는 이 값이 약 $2.5\times10^9$까지 커져 signed \code{int} 범위를 넘을 수 있다.

안전한 형태의 한 예는 다음과 같다.
\begin{lstlisting}
s += 1LL * (i + 1) * (N - i) * A;
\end{lstlisting}

\subsection{학습 질의응답: cast는 어디에 붙여야 하는가?}
\begin{lstlisting}
static_cast<long long>(a + b) * c
\end{lstlisting}
는 \code{a+b}가 먼저 \code{int}로 계산된 뒤 cast된다. 합 자체가 overflow할 수 있다면
\begin{lstlisting}
(static_cast<long long>(a) + b) * c
\end{lstlisting}
처럼 더 일찍 승격해야 한다.

\subsection{E-3. Equal Pair Score}
\begin{problembox}{중간평가 재평가 문제 E-3 --- Equal Pair Score}
정렬되지 않은 수열 $A$가 주어진다. 같은 값을 가진 모든 index pair $i<j$에 대해 그 공통값 $A_i$를 score에 더한다. 전체 score를 출력하라.

\textbf{구현 및 제출 조건}
\begin{itemize}
  \item C++17 또는 C++20을 사용한다.
  \item 충분히 빠른 알고리즘을 사용한다.
  \item 최종값뿐 아니라 실제 구현의 raw intermediate product의 최대 크기도 설명한다.
\end{itemize}

\textbf{제약 조건}
\begin{TextBlock}
1 <= N <= 200000
0 <= A_i <= 100000000
\end{TextBlock}
\end{problembox}

\begin{answerbox}[학습자 제출 답안 - 보존된 핵심]
학습자는 정렬 후 같은 값을 하나의 group으로 묶어, 값이 \code{temp}이고 개수가 \code{count}인 group마다 다음 contribution을 더했다.
\begin{lstlisting}
sum += temp * (count - 1) * count / 2;
\end{lstlisting}
관련 변수는 \code{long long}으로 두었다. 현재 기록에는 전체 제출 코드가 완전하게 남아 있지 않아 확인 가능한 핵심만 수록한다.
\end{answerbox}

\subsection{보강: 나누기 전 raw product를 확인하기}
최종값의 보수적 최대는
\[
10^8\binom{200000}{2}=1.99999\times10^{18}
\]
이지만, 실제 코드에서는 나누기 2 전에
\[
10^8\cdot199999\cdot200000
\approx3.99998\times10^{18}
\]
가 만들어진다. 최종값만 확인하면 이 intermediate를 놓칠 수 있다.

\subsection{Remediation: numeric-safety proof를 네 단계로 고정하기}
이후에는 다음 순서를 고정해 연습했다.
\begin{enumerate}
  \item constraint
  \item final bound
  \item cumulative intermediate bound
  \item actual C++ type와 capacity 비교
\end{enumerate}

예를 들어
\begin{lstlisting}
long long x = 1LL * n * (n + 1) * x_value / 2;
\end{lstlisting}
에서는 최종값뿐 아니라 \code{1LL*n}, 그 결과와 \code{n+1}의 곱, 다시 \code{x\_value}를 곱한 raw product가 각각 어느 범위에 있는지 확인한다.

\subsection{E-4. Distance Pair Score}
\begin{problembox}{중간평가 재평가 문제 E-4 --- Distance Pair Score}
길이 $N$의 정수 수열 $A$가 주어진다. 모든 $i<j$에 대해 pair score를
\[
(j-i)(A_i+A_j)
\]
로 정의한다. 모든 pair score의 합을 출력하라.

\textbf{구현 및 제출 조건}
\begin{itemize}
  \item C++17 또는 C++20을 사용한다.
  \item 시간복잡도는 $O(N)$이어야 한다.
  \item 정확한 정수값을 출력한다.
  \item numeric-safety proof에는 전역 answer upper bound, 실제 코드의 overflow-risk intermediate, 각 maximum magnitude, 실제 C++ type, Safe/Unsafe 결론을 포함한다.
  \item total space와 auxiliary space를 구분한다.
\end{itemize}

\textbf{제약 조건}
\begin{TextBlock}
1 <= N <= 50000
0 <= A_i <= 10000
\end{TextBlock}

\textbf{예제}
\begin{TextBlock}
Input
4
2 1 3 2

Output
40
\end{TextBlock}
\end{problembox}

\begin{answerbox}[학습자 제출 답안]
\begin{lstlisting}
#include <iostream>
#include <string>
#include <algorithm>
#include <vector>
using namespace std;

int main()
{
    int N = 0;
    cin >> N;
    long long sum1 = 0;
    long long sum2 = 0;
    long long temp = 0;
    long long ans = 0;
    for (long long j = 0; j < N; j += 1) {
        long long A;
        cin >> A;
        if (j == 0) {
            temp = A;
        }
        else {
            sum1 += temp;
            sum2 += (j - 1) * temp;
            temp = A;
        }
        ans += j * sum1 + j * j * A
             - sum2 - j * (j - 1) * A / 2;
    }
    cout << ans;
}
\end{lstlisting}

제출한 복잡도 분석은 시간 $O(N)$, total/auxiliary space $O(1)$이었다.
\end{answerbox}

\subsection{보강: ``대략''도 upper bound보다 작으면 안 된다}
최대 index 부근에서 실제 중요한 곱은 대략 다음 크기까지 갈 수 있다.
\[
j\cdot\code{sum1}\lesssim2.5\times10^{13},
\]
\[
j^2A\lesssim2.5\times10^{13},
\]
\[
j(j-1)A\lesssim2.5\times10^{13}.
\]
이 값들을 $5\times10^{12}$처럼 실제 최대보다 작게 적으면 유효한 upper bound가 아니다.

모든 $A_i=10000$일 때의 실제 최종값은
\[
416666666500000000\approx4.17\times10^{17}
\]
이다.

\subsection{Deeper Remediation: upper bound 검산}
큰 수의 곱은 계수와 지수를 분리해서 검산한다.
\[
(2\times10^5)^2(5\times10^4)
=4\times10^{10}\cdot5\times10^4
=2\times10^{15}.
\]

또한 ``정확한 최댓값''이 아니더라도 실제 값을 덮는 보수적 upper bound면 충분하지만, upper bound라는 이름을 붙이려면 실제 가능한 값보다 작아서는 안 된다.


\subsection{보강 질의응답 기록: cumulative intermediate와 evaluation type}
학습 과정에서 numeric-safety proof를 안정화하기 위해 다음과 같은 짧은 문답을 반복했다.

\paragraph{연습 1. 세제곱에 가까운 곱의 cumulative bound}
다음 식을 생각한다.
\begin{lstlisting}
1LL * N * (N + 1) * (2 * N + 1) / 6
\end{lstlisting}
각 factor의 최대값만 따로 나열하는 것으로 끝내지 않고, 왼쪽부터 실제로 만들어지는 cumulative product의 크기를 확인해야 한다. 핵심은 ``factor 하나가 안전하다''와 ``그 factor들을 곱한 intermediate가 안전하다''를 구분하는 것이다.

\paragraph{연습 2. 나누기 전 raw product}
조건이
\[
n\le300000,\qquad x\le5\times10^7
\]
이고
\begin{lstlisting}
1LL * x * n * (n - 1) / 2
\end{lstlisting}
를 계산한다면 최종값은 약 $2.25\times10^{18}$ 규모지만, 나누기 전 raw product는 약 $4.5\times10^{18}$ 규모다. 둘 다 signed \code{long long} 안에 들어가지만, proof에는 raw product도 포함해야 한다.

\paragraph{연습 3. 너무 늦은 cast}
\begin{lstlisting}
static_cast<long long>(a + b) * c
\end{lstlisting}
에서 $a,b\le1.5\times10^9$라면 \code{a+b} 자체가 signed \code{int} 범위를 넘을 수 있다. cast는 그 뒤에 적용되므로 충분히 이르지 않다. 학습자는 이 경우 \code{a} 또는 \code{b}를 addition 전에 cast해야 함을 확인했다.

\paragraph{연습 4. 타입은 \code{long long}이어도 magnitude가 너무 클 수 있다}
조건이
\[
n\le400000,\qquad x\le10^8
\]
일 때
\begin{lstlisting}
1LL * n * (n + 1) * x / 2
\end{lstlisting}
의 최종값은 약 $8\times10^{18}$ 규모지만, 나누기 전 raw product는 약 $1.6\times10^{19}$까지 갈 수 있다. 선두 \code{1LL}로 타입을 승격해도 magnitude가 signed 64-bit capacity를 넘으면 unsafe하다.

\paragraph{연습 5. 먼저 약분하여 intermediate 줄이기}
같은 구조를 다음처럼 바꿀 수 있다.
\begin{lstlisting}
long long a = n;
long long b = n + 1;
if (a % 2 == 0) a /= 2;
else b /= 2;
long long ans = a * b * x;
\end{lstlisting}
먼저 짝수 factor를 2로 나눈 뒤 곱하면 최대 intermediate를 최종값 규모인 약 $8\times10^{18}$로 낮출 수 있다.

\subsection{보강 질의응답 기록: upper bound 산술 체크포인트}
다음 문답은 ``내가 적은 값이 실제 upper bound인가''를 빠르게 검산하기 위해 사용했다.

\paragraph{A. 두 정수의 곱}
\begin{lstlisting}
long long x = a * b;
\end{lstlisting}
\[
0\le a\le300000,\qquad0\le b\le700000.
\]
$2.1\times10^{11}$은 유효한 upper bound인가? 학습자는 ``예''라고 답했고,
\[
(3\times10^5)(7\times10^5)=2.1\times10^{11}
\]
로 확인했다.

\paragraph{B. 제곱이 포함된 곱}
\begin{lstlisting}
long long x = a * a * b;
\end{lstlisting}
\[
0\le a\le200000,\qquad0\le b\le50000.
\]
학습자는
\[
(2\times10^5)^2(5\times10^4)=2\times10^{15}
\]
로 계산해 $2\times10^{15}$이 유효한 upper bound임을 확인했다.

\paragraph{C. 느슨한 upper bound도 가능한가?}
\[
p<60000,\qquad q<40000,\qquad r\le100000.
\]
\[
pqr<2.4\times10^{14}.
\]
따라서 $2.4\times10^{14}$도, 더 느슨한 $3\times10^{14}$도 모두 유효한 upper bound다.

\paragraph{D. 과학적 표기법 정규화}
\[
i<100000,\qquad A\le10^6
\]
일 때
\[
i(i+1)A<10^{16}
\]
로 잡을 수 있다. $10^{16}$의 과학적 표기법은 $1\times10^{16}$이다. 실제 정수 최대를 사용하면
\[
99999\cdot100000\cdot10^6=9.9999\times10^{15}
\]
이다.

\subsection{보강 질의응답 기록: bound와 C++ 타입을 함께 확인하기}
\paragraph{E. destination만 \code{long long}}
\begin{lstlisting}
int a = 100000;
int b = 100000;
long long x = a * b;
\end{lstlisting}
학습자는 최대 intermediate가 $10^{10}$이고 실제 곱셈 타입은 \code{int}이므로 Unsafe라고 답했다.

\paragraph{F. 선두 \code{1LL}}
\begin{lstlisting}
int a = 100000;
int b = 100000;
long long x = 1LL * a * b;
\end{lstlisting}
학습자는 최대값 $10^{10}$, multiplication chain의 타입 \code{long long}, Safe라고 답했다.

\paragraph{G. 괄호 내부 타입과 바깥 multiplication chain}
\begin{lstlisting}
int n = 100000;
long long x = 1LL * n * (n + 1) * (2 * n + 1);
\end{lstlisting}
학습자는 보수적으로 $6\times10^{15}$을 upper bound로 잡았다. \code{n+1}과 \code{2*n+1}은 \code{int}로 계산되지만 최대 $200001$이라 안전하고, 선두 \code{1LL} 이후 곱셈 체인은 \code{long long}으로 진행된다.

\subsection{E-5. Weighted Prefix Load}
\begin{problembox}{중간평가 재평가 문제 E-5 --- Weighted Prefix Load}
길이 $N$의 음이 아닌 정수 수열 $A_0,\dots,A_{N-1}$이 주어진다. 각 위치 $i$에 대해
\[
P_i=A_0+\cdots+A_i
\]
라고 정의한다. 다음 값을 계산하라.
\[
S=\sum_{i=0}^{N-1}(i+1)P_i.
\]

\textbf{구현 및 제출 조건}
\begin{itemize}
  \item C++17 또는 C++20을 사용한다.
  \item 시간복잡도는 $O(N)$이어야 한다.
  \item auxiliary space는 $O(1)$이어야 한다.
  \item numeric-safety proof에는 전역 answer upper bound, 주요 intermediate upper bound, 실제 C++ type, capacity 비교를 포함한다.
  \item total space와 auxiliary space를 구분한다.
\end{itemize}

\textbf{제약 조건}
\begin{TextBlock}
1 <= N <= 200000
0 <= A_i <= 100
\end{TextBlock}

\textbf{예제 입력}
\begin{TextBlock}
3
2 1 3
\end{TextBlock}

\textbf{예제 출력}
\begin{TextBlock}
26
\end{TextBlock}
\end{problembox}

\begin{answerbox}[학습자 제출 답안]
\begin{lstlisting}
#include <iostream>
#include <string>
#include <algorithm>
#include <vector>

using namespace std;

int main()
{
    int N = 0;
    cin >> N;
    long long P = 0;
    long long S = 0;
    for (long long i = 0; i < N; i += 1) {
        long long A;
        cin >> A;
        P += A;
        S += (i + 1) * P;
    }
    cout << S;
}
\end{lstlisting}

\textbf{시간복잡도}: $O(N)$.\\
\textbf{공간복잡도}: total $O(1)$, auxiliary $O(1)$.\\
\textbf{numeric-safety 설명의 핵심}: $P\le2\times10^7$, $(i+1)P\le4\times10^{12}$, 그리고 최대 $2\times10^5$개 항을 합쳐도 보수적으로 $S\le8\times10^{17}$이므로 signed 64-bit 범위 안이다.\\
\textbf{edge case}: $N=1$, $A_0=11$이면 출력은 11.
\end{answerbox}

\subsection{질의응답: 모든 operand가 \code{long long}이어야 하는가?}
그럴 필요는 없다. 예를 들어 \code{i}가 \code{int}여도 \code{i+1} 자체가 \code{int} 범위 안에 있고, 그 다음 \code{long long P}와 곱해질 때 승격된다면 안전할 수 있다. 중요한 것은 ``모든 변수를 큰 타입으로 만드는 것''이 아니라 \textbf{각 subexpression이 어떤 타입으로 언제 계산되는지}를 확인하는 것이다.

% =========================================================
% Part F
% =========================================================
\section{Part F --- 최종 성취도 평가 문제}

\subsection{F-1. Grouped Weighted Score}
\begin{problembox}{최종평가 문제 F-1 --- Grouped Weighted Score}
$N$개의 그룹이 주어진다. $g$번째 그룹에는 $M_g$개의 정수 $A_{g,0},\dots,A_{g,M_g-1}$이 있다. 다음 값을 계산하라.
\[
S=\sum_{g=0}^{N-1}\sum_{j=0}^{M_g-1}(g+1)(j+1)A_{g,j}.
\]

\textbf{구현 및 제출 조건}
\begin{itemize}
  \item C++17 또는 C++20을 사용한다.
  \item 전체 원소 수를 $T=\sum_g M_g$라 할 때 시간복잡도는 $O(N+T)$ 이하여야 한다.
  \item 정확한 정수값을 출력한다.
  \item numeric-safety proof에서 전역 answer upper bound, 실제 코드의 overflow-risk intermediate, 각 upper bound, actual C++ type, Safe/Unsafe 결론을 설명한다.
  \item total space와 auxiliary space를 구분한다.
\end{itemize}

\textbf{제약 조건}
\begin{TextBlock}
1 <= N <= 200000
1 <= M_g
T = sum M_g <= 200000
0 <= A_g,j <= 100
\end{TextBlock}

\textbf{예제 입력}
\begin{TextBlock}
3
2 4 1
1 3
3 2 0 5
\end{TextBlock}

\textbf{예제 출력}
\begin{TextBlock}
63
\end{TextBlock}
\end{problembox}

\begin{answerbox}[학습자 제출 답안]
\begin{lstlisting}
#include <iostream>
#include <string>
#include <algorithm>
#include <vector>

using namespace std;

int main()
{
    int N = 0;
    cin >> N;
    long long S = 0;
    for (int g = 0; g < N; g += 1) {
        int Mg = 0;
        cin >> Mg;
        for (int j = 0; j < Mg; j += 1) {
            int A = 0;
            cin >> A;
            S += 1LL * (g + 1) * (j + 1) * A;
        }
    }
    cout << S;
}
\end{lstlisting}

\textbf{시간복잡도}: 내부 body는 전체 원소마다 한 번 실행되므로 $O(T)$이며, 모든 그룹이 비어 있지 않으므로 $N\le T$라서 요구한 $O(N+T)$ 안에 들어간다.\\
\textbf{공간복잡도}: 입력을 저장하지 않으므로 total/auxiliary 모두 $O(1)$.\\
\textbf{numeric-safety 설명의 핵심}: 한 항은 보수적으로 $(2\times10^5)(2\times10^5)(100)=4\times10^{12}$ 이하이고, 항 수가 최대 $2\times10^5$이므로 $S\le8\times10^{17}$. 선두의 \code{1LL} 이후 multiplication chain은 64-bit 범위에서 안전하다.\\
\textbf{edge case}: $N=1$, $M_0=1$, $A_{0,0}=61$이면 출력은 61.
\end{answerbox}

\subsection{최종평가 이후 정리: 충분한 타입과 필요한 타입을 구분하기}
이 문제에서는 \code{g}, \code{j}, \code{A}를 모두 \code{int}로 두면서 선두의 \code{1LL}로 위험한 곱셈 체인만 승격했다. 이는 ``모든 변수를 \code{long long}으로 만들기''보다 실제 expression을 기준으로 타입을 선택한 예다.

% =========================================================
% Part G
% =========================================================
\section{Part G --- Adaptive Extra Problems}

\subsection{G-1. Reverse Archive Score}
\begin{problembox}{Adaptive Extra Problem G-1 --- Reverse Archive Score}
$N$개의 비어 있지 않은 소문자 문자열 $s_0,\dots,s_{N-1}$이 주어진다. 먼저
\[
S=\sum_{i=0}^{N-1}(i+1)|s_i|
\]
를 계산한다. 그 후 각 문자열의 문자 순서를 뒤집는다. 첫 줄에 $S$를 출력하고, 이후 문자열의 입력 순서도 역순으로 하여 뒤집힌 문자열들을 한 줄에 하나씩 출력하라.

\textbf{구현 및 제출 조건}
\begin{itemize}
  \item C++17 또는 C++20을 사용한다.
  \item 전체 문자 수를 $L=\sum_i|s_i|$라 할 때 시간복잡도는 $O(N+L)$이어야 한다.
  \item 모든 문자열을 저장하는 것은 허용한다.
  \item \code{reverse}를 포함한 주요 STL 연산의 비용을 실제 처리 원소 수 기준으로 설명한다.
  \item numeric-safety proof에서 $S$의 전역 upper bound와 실제 계산 타입을 설명한다.
  \item total space와 auxiliary space를 구분한다.
\end{itemize}

\textbf{제약 조건}
\begin{TextBlock}
1 <= N <= 200000
All strings are nonempty lowercase strings.
L = total number of characters <= 200000
\end{TextBlock}

\textbf{예제 입력}
\begin{TextBlock}
3
abc
z
xy
\end{TextBlock}

\textbf{예제 출력}
\begin{TextBlock}
11
yx
z
cba
\end{TextBlock}
\end{problembox}

\begin{answerbox}[학습자 제출 답안]
\begin{lstlisting}
#include <iostream>
#include <string>
#include <algorithm>
#include <vector>

using namespace std;

int main()
{
    int N = 0;
    cin >> N;
    long long S = 0;
    vector<string> v = {};
    for (int i = 0; i < N; i += 1) {
        string si = "s";
        cin >> si;
        S += 1LL * (1 + i) * si.size();
        reverse(si.begin(), si.end());
        v.push_back(si);
    }
    cout << S << '\n';
    for (int i = 1; i < N + 1; i += 1) {
        cout << v[N - i] << '\n';
    }
}
\end{lstlisting}

학습자는 \code{reverse}가 각 문자열 길이에 비례한다는 점을 인식하고 전체를 $L$로 합산했다. 이후 분석 보강에서 출력 비용, aggregate storage, \code{string::size\_type}을 다시 점검했다.
\end{answerbox}

\subsection{G-1 보강: aggregate size를 정확히 세기}
모든 문자열을 출력하면 총 $L$개의 문자를 출력하므로 출력 자체가 $O(L)$이다. 또한 \code{vector<string>}은 문자열 객체 $N$개와 총 $L$개의 문자를 저장하므로 aggregate storage는
\[
O(N+L)
\]
이며, 문자열이 모두 비어 있지 않아 $N\le L$이므로 이 문제에서는 $O(L)$로 단순화할 수 있다.

score는 더 직접적으로
\[
S=\sum_i(i+1)|s_i|\le N\sum_i|s_i|=NL\le4\times10^{10}
\]
로 bound할 수 있다.

\subsection{G-2. AtCoder ABC238 B --- Pizza}
\begin{problembox}{External Extra Problem G-2 --- Pizza}
원형 피자에 $0^\circ$ 방향의 기준 칼집이 있다. $A_1,\dots,A_N$만큼 차례대로 피자를 회전한 뒤 매번 같은 고정 방향에 칼집을 낸다. 모든 칼집이 만들어진 뒤 가장 큰 피자 조각의 중심각을 출력하라.

\textbf{구현 및 제출 조건}
\begin{itemize}
  \item C++17 또는 C++20을 사용한다.
  \item 시간복잡도와 공간복잡도를 설명한다.
  \item 모든 정수 계산이 \code{int} 또는 \code{long long} 중 어느 범위에서 안전한지 설명한다.
  \item edge case를 직접 검증한다.
\end{itemize}

\textbf{제약 조건}
\begin{TextBlock}
1 <= N <= 359
1 <= A_i <= 359
No two cuts have the same position.
\end{TextBlock}

출처: AtCoder ABC238 B --- Pizza.
\end{problembox}

\begin{answerbox}[학습자 제출 답안]
\begin{lstlisting}
#include <iostream>
#include <string>
#include <algorithm>
#include <vector>

using namespace std;

int main()
{
    int N = 0;
    cin >> N;
    vector<int> v = {0};
    int A = 0;
    int m = 0;
    for (int i = 1; i < N + 1; i += 1) {
        int Ai = 0;
        cin >> Ai;
        if (A + Ai >= 360) {
            A = A + Ai - 360;
        }
        else {
            A += Ai;
        }
        v.push_back(A);
    }
    v.push_back(360);
    sort(v.begin(), v.end());
    for (int i = 1; i < v.size() - 1; i += 1) {
        int v1 = v[i - 1];
        int v2 = v[i];
        int v3 = v[i + 1];
        int a1 = v2 - v1;
        int a2 = v3 - v2;
        if (m < max(a1, a2)) {
            m = max(a1, a2);
        }
    }
    cout << m;
}
\end{lstlisting}

\textbf{시간복잡도}: 정렬이 지배하므로 $O(N\log N)$.\\
\textbf{공간복잡도}: cut position vector가 $O(N)$.\\
\textbf{정수 범위}: cumulative angle을 매번 $0$ 이상 $359$ 이하로 유지하고, 인접 차이도 최대 $360$이므로 모두 \code{int}에 충분히 들어간다.\\
\textbf{edge case}: $N=1$, $A_1=20$이면 가장 큰 조각은 $340^\circ$.
\end{answerbox}

\subsection{G-2 보강: 반복 횟수와 asymptotic 결론}
\code{v}에는 $0$, $N$개의 cut position, $360$이 들어가므로 크기는 $N+2$다. 두 번째 scan은 $N$번 실행된다. 이를 $N-2$라고 잘못 세더라도 $O(N)$이라는 asymptotic 결론은 같지만, 정확한 iteration count를 세는 습관은 유지해야 한다.

또한 \code{i < v.size()-1}은 signed \code{int}와 unsigned size type을 비교하므로 compiler warning이 날 수 있다. 이 문제의 작은 범위에서는 correctness 문제가 없지만, 일반화할 때는 타입을 의식한다.

% =========================================================
% Part H
% =========================================================
\section{Part H --- 숙달 점검과 복습}

\subsection{H1. 개념 자기점검}
다음 질문에 자료 없이 답할 수 있는지 확인하라.
\begin{enumerate}
  \item $O(N)$과 $O(N^2)$의 차이를 반복 횟수 관점에서 설명할 수 있는가?
  \item 순차로 실행되는 두 $O(N)$ loop가 왜 $O(N)$인가?
  \item 그룹별 크기 $M_g$가 있을 때 총 iteration 수를 $T=\sum M_g$로 표현할 수 있는가?
  \item 문자열 $N$개 전체 문자 수가 $L$일 때 모든 문자열을 뒤집는 비용은 왜 $O(L)$인가?
  \item 모든 문자열을 출력하는 비용을 왜 단순히 $O(N)$이라고 하면 안 되는가?
  \item total space와 auxiliary space의 차이는 무엇인가?
  \item \code{long long x = a * b;}가 unsafe할 수 있는 이유는 무엇인가?
  \item \code{1LL * a * b}가 어떤 점에서 다른가?
  \item cast가 연산 뒤에 적용되면 왜 너무 늦을 수 있는가?
  \item 최종값이 안전해도 intermediate가 overflow할 수 있는 예를 설명할 수 있는가?
  \item 가능한 입력 한 개와 global upper bound는 어떻게 다른가?
  \item 실제 최대보다 작은 값은 왜 upper bound가 아닌가?
  \item \code{string::size()}의 반환형을 단순한 \code{int}라고 보면 왜 위험한가?
  \item signed/unsigned arithmetic conversion을 언제 의식해야 하는가?
\end{enumerate}

\subsection{H2. 구현 자기점검}
\begin{checklistbox}[S0-C 코드 작성 후 체크리스트]
\begin{itemize}
  \item[$\square$] 문제의 실제 input-size 변수를 모두 정의했는가? ($N$, $L$, $T$ 등)
  \item[$\square$] 반복문이 외형상 중첩됐다는 이유만으로 $O(N^2)$라고 단정하지 않았는가?
  \item[$\square$] 문자열 복사/출력/reverse 비용을 payload 길이까지 포함했는가?
  \item[$\square$] 정렬 비용 $O(N\log N)$을 포함했는가?
  \item[$\square$] total space와 auxiliary space를 구분했는가?
  \item[$\square$] 최종 answer의 전역 upper bound를 적었는가?
  \item[$\square$] 실제 코드의 위험한 intermediate expression을 적었는가?
  \item[$\square$] multiplication/division 순서를 포함한 cumulative bound를 확인했는가?
  \item[$\square$] destination type이 아니라 actual expression type을 확인했는가?
  \item[$\square$] 필요한 경우 \code{1LL} 또는 cast를 충분히 이른 위치에 넣었는가?
  \item[$\square$] \code{size()}가 반환하는 unsigned size type을 확인했는가?
  \item[$\square$] edge case가 실제 문제의 입력 조건을 만족하는가?
\end{itemize}
\end{checklistbox}

\subsection{H3. S0-C에서 형성해야 할 핵심 사고}
S0-C의 목표는 ``큰 값이면 \code{long long}을 쓴다''를 암기하는 것이 아니다. 다음 순서를 습관으로 만드는 것이 핵심이다.
\[
\boxed{\text{input size} \to \text{operation count} \to \text{numeric magnitude} \to \text{actual C++ type}}
\]

복잡도와 numeric safety는 모두 ``코드가 실제로 무엇을 처리하고, 어떤 순서로 평가되는가''를 정확히 세는 훈련이라는 공통점을 가진다.

\subsection{H4. 다음 Unit과의 연결}
다음 Learning Unit인 S1-A에서는 associative container를 다룬다. 다음 도구들의 선택에는 operation complexity, ordering, 중복 처리 조건이 직접 연결된다.
\begin{itemize}
  \item \code{unordered\_map}, \code{unordered\_set}
  \item \code{map}, \code{set}, \code{multiset}
\end{itemize}

S0-C에서 형성한 다음 질문을 그대로 가져가야 한다.
\begin{itemize}
  \item lookup/update가 몇 번 필요한가?
  \item 평균 $O(1)$과 $O(\log N)$ 중 무엇이 필요한가?
  \item ordering이 필요한가?
  \item container에 저장되는 총 데이터 크기는 얼마인가?
  \item key/value 계산에서 정수 범위와 타입은 안전한가?
\end{itemize}
